\documentclass[a4paper,12pt]{ctexart}
\usepackage{xcolor}
\usepackage{graphicx}
\usepackage[top=1.8cm, bottom=1.8cm, left=1.8cm, right=1.8cm]{geometry}
\usepackage{array}
\usepackage{hyperref}
\hypersetup{colorlinks=true,linkcolor=blue!70!black}
\linespread{1.35}

\newcommand{\ruby}[2]{#1\textsuperscript{#2}}
\newcommand{\xuehai}[2]{%
  \vspace{0.35cm}
  \noindent
  \fcolorbox{blue!50!black}{blue!5}{%
    \begin{minipage}{0.915\textwidth}
    \textbf{\textcolor{blue!60!black}{学海拾遗：}#1} \\[0.1cm]
    {\small #2}
    \end{minipage}%
  }
  \vspace{0.35cm}
}

\title{\textcolor{blue!70!black}{\Huge\textbf{道德与法治}}}
\author{}
\date{}

\begin{document}
\maketitle
\thispagestyle{empty}

\tableofcontents
\newpage

\section{\textcolor{orange!80!black}{第一课\quad 规则——没有规则的世界长什么样？}}

\subsection*{操场上的混乱}

课间打了铃。四个班的同学同时涌向了操场上的两个秋千架。二班的浩浩跑了第一，一屁股坐在秋千上。三班的婷婷慢了一步——但她拽住了另一根秋千的链子。还有七八个同学在后面排队。

如果没有任何规则——浩浩可以一直荡到上课铃再响（因为他到得最早）。婷婷可以把链子一直拽在手里不让任何后面的人用（因为"我拿到的就是我的"）。但这样一来——排队的同学就会愤怒。没有规则的世界——不是自由的——是紊乱的。力气最大的、最先冲到的、声音最响的人拿走一切。\textbf{规则保护了那些爆发力不够强——但同样有权利的人。}

\subsection*{家庭里的规则、班级里的规则}

你爸妈在家里也定了一些规则：晚上几点不玩手机、周末先写完作业再出去。为什么？不是因为他们"管你管得不耐烦了"。因为如果没有这些规则——一晚上不睡第二天上课打瞌睡，周末玩到天黑然后冲刺写作业写到半夜。\textbf{规则在我们管不住自己的时候替我们管。}

班级也有规则——值日表轮流、上课举手再发言。班规的作用和操场上秋千排队是一样的——\textbf{不能同时使用同一件东西的时候，规则来排序。}

\subsection*{交通规则——最需要遵守的规则}

放学走到十字路口。斑马线对面的红灯亮了。一位外卖骑手想"反正没车"——闯了红灯，差点撞翻一个低年级的小女孩。交通规则和其他规则有一个巨大的不同：\textbf{一旦违反了交通规则——后果不是"被老师骂几句"就完了。可能有人流血。}交叉路口对红灯的共识——所有车辆一起停车、所有行人一起通行——是一个社会级的\textbf{互不伤害协议}。

\subsection*{一块刻在广场上的石板——规则为什么要公开？}

公元前五世纪的罗马。贵族垄断了所有法律的解释权——因为法律条文没有被写下来，全凭他们背书背诵。平民去打官司，贵族说"根据古老的法律——你败诉了"，但这条"古老的法律"具体写的是什么，没有任何人能看到原文。平民觉得不对——但他们没法反驳，因为他们不知道法条上的字到底长什么样。

忍无可忍的罗马平民展开了一场漫长的抗议——他们集体离开城市，拒绝为贵族打仗。最终贵族被迫让步。公元前450年左右，罗马人在市中心广场上立起了十二块青铜板——把所有的法律条文\textbf{一条一条刻在上面，任何人都可以走过去站在旁边读}。这就是著名的《\textbf{十二铜表法}》——古罗马最早的成文法典。

\textbf{把规则刻在所有人看得见的地方——一个看似简单的动作，把规则从"少数人的秘密武器"变成了"任何人的盾牌"。}

\xuehai{从习惯到成文法——规则的三种形态}{
  人类社会的规则大致经历了三个演化阶段。最初的规则是\textbf{习惯}——没有人刻意制定，但大家都默认这么做（比如排队时后来的站在后面）。习惯的缺点是模糊——"站在后面"是指两米还是一厘米？

  第二阶段是\textbf{不成文的传统法律}——长老或祭司背诵的古老规矩，只有少数人能"解释"。古罗马贵族垄断法律解释权就是这样的状态。受规则约束的人看不到规则原文，这是不公平的最深层来源。

  第三阶段是\textbf{成文法}——规则被用文字写下来、公之于众。它带来的变化不是"规则变多了"，而是\textbf{规则可以被任何识字的人独立查阅}。统治者不能再随便说"古老的法律规定你必须……"——因为现在被统治的人可以自己走过去读那条石板上的原文。
}

\subsection*{想一想}

\begin{enumerate}
  \item 罗马平民为什么要逼贵族把法条刻在广场石板上？把一条规则"写下来并且贴出去"——它的力量到底在哪？
  \item 你们班如果有一条班规从来没有人贴出来过——只有班长一个人"知道"——其他同学遵守的时候心里舒服吗？
  \item 你们学校的校规贴在哪面墙上？你上一次站在它前面读完一条完整的是什么时候？如果没有任何人读过但仍然在执行——它和"不成文法"之间的距离是什么？
\end{enumerate}

\newpage
\section{\textcolor{orange!80!black}{第二课\quad 对与错——这些判断从哪里来？}}

\subsection*{偷吃同学的零食}

体育课，全班在跑圈。你的同桌把一包刚买的小鱼干放在课桌抽屉里。你经过他的座位——看见了那包零食。你伸手拿了一块塞进嘴里。这件事为什么是错的？

第一种答案：\textbf{因为后果。}他会生气——以后不信任你了——你伤了他的心。做这件事带来了坏的后果——所以错。第二种答案：\textbf{因为这件事本身就错。}不问别人就拿别人的东西——不管他发没发现、生不生气——偷拿本身就是错的。

\subsection*{朋友偷了商店的巧克力——你说不说？}

你和好朋友放学一起逛便利店。你已经买完东西在门口等他。他从店里出来的时候，口袋里鼓了一小块——是一块没有付钱的巧克力。你小声问他："你付钱了吗？"他摇摇头："没事，没人看到。"

现在你站在路口。如果你不告诉任何人——朋友明天还会继续。如果你告诉老师或他父母——朋友可能会跟你绝交。\textbf{忠诚和诚实在你胸口撞在了一起。}哪一种"对"更有分量？

\xuehai{道德判断的六个台阶——科尔伯格的实验}{
  美国心理学家劳伦斯·科尔伯格在1950年代给不同年龄的孩子讲道德困境（类似"朋友偷巧克力"这种问题），然后记录他们给出的理由。他发现人的道德判断大致沿着六个台阶往上走：

  前三个阶段是"因为怕挨罚"（第一阶段：不偷——因为老师会罚站）→"因为对我有好处"（第二阶段）→"因为别人会夸我是好孩子"（第三阶段：不偷——因为老师会觉得我诚实）。后三个阶段逐渐脱离自我中心：开始考虑"规则是用来维护秩序的"（第四阶段）、"法律是人定的——不合理的法律可以被质疑"（第五阶段）、"有些原则——比如人的尊严——高于任何具体的法律条文"（第六阶段）。

  科尔伯格的数据显示，10岁之前多数孩子的判断处在第一到第三阶段之间。12-14岁开始有一部分孩子进入第四阶段——\textbf{不再因为"被老师抓到"才判断对错，而是因为内心已经装了一个"法律维护秩序"的自主开关}。
}

\subsection*{多给了你8分——你没吭声}

期中测试。数学老师发卷子的时候错把你的68分算成了76分——在总分一栏多写了8分。你当场发现了。但是76分刚好跨过了你妈说的"考不到70分别想出去玩"的那条线。你没有举手纠正。

你没有撒谎——试卷上的红色分数是老师自己写上去的。你也没有偷分数——是分数"找上你的"。\textbf{但是什么都不做——算不算一种"错"？}

\subsection*{想一想}

\begin{enumerate}
  \item 朋友偷巧克力——你告诉老师，算不算"告密"？"告密"和"保护别人不继续犯错"中间那条线在哪？
  \item 监考老师算错了分你没说——从"后果"看是"好事"吗？从"行为本身"看呢？
  \item 科尔伯格说大多数10岁以下的孩子因为"怕挨罚"才做好事。你觉得自己现在到了第几阶段？有什么事情你坚持做不是因为怕被罚——而是因为你觉得它\textbf{本身就是对的}？
\end{enumerate}

\newpage
\section{\textcolor{orange!80!black}{第三课\quad 权利——你有什么权利？}}

\subsection*{一份跟你有关的国际文件}

1989年，联合国通过了一份《\textbf{儿童权利公约}》。全世界几乎所有国家都签了。它不是"大人在给你恩惠"。它是一份\textbf{有约束力的国际法文件}：你作为未满十八周岁的人，有一些基本权利不是任何人可以随便拿走的。

挑三个最直接的说：\textbf{受教育的权利}——没人可以因为你住在哪里、是否残疾、家境困难而把你关在教室门外。\textbf{表达意见的权利}——家庭、学校、社区决定任何涉及你的重大事情时，\textbf{你有权利被认真地听}。\textbf{被保护免于一切形式的暴力和剥削}——不准童工剥削，不准从事会让你失去上学机会、有危险、对身体有害的工作。

还有一条最容易被大人忽略的条款：\textbf{儿童有休息和玩耍的权利。}这不是"玩"而已，是公约第31条明文写的。当你每天被补习班排满完全没有放空的时间——公约在保护你。

\subsection*{一个巴基斯坦女孩的故事}

2012年10月9日，巴基斯坦斯瓦特山谷。一辆校车像往常一样载着一群女孩子放学回家。两个持枪的男人拦停了校车。他们走上车，大声问："\textbf{谁是马拉拉？}"一个十五岁的女孩站了起来。

他们朝她的头部开枪。

马拉拉·优素福扎伊在巴基斯坦西北部的山区长大。她十一岁那年开始用笔名在BBC的乌尔都语网站上写博客——记录在塔利班控制下，女孩子是怎样被禁止上学的。她写道："他们可以把我们的学校烧掉，但他们不能把我们的想法烧掉。"她继续上学、继续写、继续公开演讲——要求所有女孩和男孩一样有权坐在教室里。

子弹没有杀死她。她被紧急送到英国伯明翰的医院抢救，活了下来。2014年，十七岁的马拉拉获得了诺贝尔和平奖——她是这个奖项最年轻的获得者。联合国把她的生日7月12日定为"\textbf{马拉拉日}"。

\subsection*{权利的背面——义务}

你有每天吃午饭的权利。但这顿午饭不是凭空从餐厅里冒出来的——厨工起了大早来备菜、碗筷是校工一桌一桌码的。\textbf{你占用这一顿饭，同时也有义务：安静排队打饭、不浪费食物。}

\subsection*{权利不能被投票投掉}

一个班百分之八十的同学\textbf{投票决定}新转学来的张晓晓没有课间去操场的权利——这票通过有效吗？无效。因为张晓晓不是一个用大伙投票多寡来决定拥不拥有的东西。\textbf{有些权利——不管你代表多少个人投票反对——它仍然存在。}

\xuehai{消极权利和积极权利——两种不同逻辑的保护}{
  法律哲学中把基本权利分成了两类：

  \textbf{消极权利}——别人和政府\textbf{不能对你做什么}。不被非法逮捕、不被虐待、不被任意剥夺财产。这些权利的运行方式是一个"停止"按钮——只要没有人踩进来，权利就是安全的。

  \textbf{积极权利}——别人或政府\textbf{必须为你提供什么}。受教育权需要学校、教材和老师；医疗权需要医院和药品。这些权利不是"停下来就行"——它需要资源被预算和分配。一个乡村学校没有英文老师——孩子没有被任何人阻挡上学——但她的受教育权已经在实然上被缩水了。

  联合国《儿童权利公约》同时包含了这两种权利——这就是为什么它既是一份"不许伤害孩子"的禁令、也是一份"大人们需要凑钱建学校和医院"的账单。
}

\subsection*{想一想}

\begin{enumerate}
  \item 马拉拉中枪之后全世界去救她——如果她没有写博客、没有上新闻、只是一个普通的巴基斯坦山村女孩——她的权利还存不存在？
  \item 成年人抽烟。在公交车和食堂里抽烟——是"他自己的身体"的自由吗？烟雾飘到旁边座位上的那个婴儿吸进了肺——所谓"一个人自己的自由"到鼻子尖之前为止？
  \item 你家附近有一个孩子因为户口问题不能在所在的城市上公立初中。没有人禁止她上学——但她的受教育权在哪个环节被卡住了？这是消极权利出问题还是积极权利出问题？
\end{enumerate}

\newpage
\section{\textcolor{orange!80!black}{第四课\quad 公平——怎么分才公平？}}

\subsection*{分蛋糕的聪明办法}

你和你妹妹（或你最好的两个朋友）——三个人分一块蛋糕。把你推出来切。你手很稳——但是其他两个人会相信你没有偏心把最大块留给自己吗？

最古老的解决办法：\textbf{切蛋糕的人最后一个选。}这套程序不需要公平的法官——它会自己推着分法走向平均。如果你切得不均，你会吃亏——\textbf{程序替道德做了裁判。}

\subsection*{一个让全世界吵了十几年的运动员}

南非运动员\textbf{卡斯特尔·塞门亚}在2009年柏林世锦赛上获得了女子800米金牌。她跑得太快了——快到了让人开始质疑\textbf{"她到底应不应该被允许参加女子比赛"}。

塞门亚天生体内的睾酮水平远高于一般的女性——这是她出生时就携带的生理条件，不是她"吃药的"。国际田联（现在的世界田联）制定了一条规则：天生睾酮过高的女性运动员如果要参加400米到一英里的女子比赛，必须服用药物把睾酮降到某个水平以下。塞门亚拒绝服药——她认为这不是"犯规"，这是我的身体。

支持田联的人说：\textbf{你不降睾酮，对其他女运动员不公平}——就像你不能让一个天生腿长到两米的人去跟一般身高的人抢篮板。支持塞门亚的人说：\textbf{强迫一个健康的女人为了"比别人弱一点"而吃药——这本身公不公平？}

这个问题已经吵了十几年，没有任何一方能完全说服另一方。\textbf{当"公平"撞上"天生差异"——那条线到底应该画在哪？}

\xuehai{分配正义的三条思路}{
  哲学家为"怎么分才算公平"这个问题吵了两千年，大致可以归类成三条路：

  \textbf{按需分配}——谁更需要就给谁多。医院急诊室的呼吸机不分你是国王还是乞丐——谁快憋死了谁先用。它的道德力量最强（"不帮最弱的人还叫什么公平"），但在非生死场景下很容易被滥用——"我缺钱所以我应该拿更多"——多么缺才算"该拿"？

  \textbf{按劳分配}——你种了几个小时的田就收几袋粮。这条逻辑保护了愿意多干的人不被"搭便车"。但它有一个盲区：你的劳动能力本身可能就不是你"挣"的——生下来智力超群的人做同样一小时的工作比智力受损的人多产十倍——这是"按劳"还是"按基因"？

  \textbf{按人头平分}——最简单也最不需要解释：一人一份。干净，但擦掉了所有差别——包括那些不能被忽略的差别（有人需要吃药才能站到起跑线上）。

  你每说一次"这不公平"，不管你有没有意识到——你都在调这三条线之间的某个刻度盘。
}

\subsection*{平等和公平是两回事}

七个小朋友一起去野餐。你带了一大袋橘子。每个人都发一个——\textbf{平等}。但如果其中的小可是一位血糖低到上午不进食就会晕倒的同学——她需要优先拿到一整个橙子和面包。你分给她多一些以满足她的基础健康——这叫\textbf{公平}。

\textbf{平等是每人数目一样}；\textbf{公平是给不同需求的人不同的补充，让每个人回到同一条起跑线上。}

\subsection*{按劳分配 vs 按人头分配}

小组做了一个月的垃圾分类科普项目。三个组员拼命努力，两个组员划过水。最后拿到的奖品是一整盒零食。怎么分？"按人头平分"——划水的人不干活照样均享团队的一切成果——对出了力气的人来说公平吗？"按劳分配"——你投入了多少时间兑换多少奖品——对偷懒的人的惩罚会不会过头了？

\subsection*{想一想}

\begin{enumerate}
  \item 塞门亚天生睾酮比别人高——她该不该为了"公平"去吃药降低它？如果强迫她吃药——这本身是不是另一种不公平？
  \item 全班投票决定用班费给这次数学考前三名的同学买奖品——投"同意"的人里有一半是这次根本没考进前三的人，他们为什么投票同意？这种"让渡给自己的利益"公平吗？
  \item 你们家的家务活按什么原则分配的？按需（谁今天有空谁做）、按劳（谁做得多谁拿奖励）还是按人头（每人三件菜一个碗）？用这三条线来诊断一下你上次觉得"这不公平"那一刻。
\end{enumerate}

\newpage
\section{\textcolor{orange!80!black}{第五课\quad 自由——我可以想做什么就做什么吗？}}

\subsection*{"我的自由"停在哪里？}

"你挥舞拳头的权利，止于我的鼻尖开始之处。"这句话被引用了无数次——因为它干净地画出了自由的边界：\textbf{你的动作有没有打到一个不想被你打到的人。}

\subsection*{你妈能不能看你手机？}

你的手机是你自己的私人空间。你妈趁你洗澡的时候翻了你的聊天记录。你发现了。你气炸了。"这是我的隐私！"你吼。

但你妈怎么说的？"你才十一岁。你在网上跟谁聊天——我有责任知道。万一你在跟一个假装是同班同学的成年人约线下见面——我什么都不知道，出事了谁来负责？"

\texbf{你的隐私权和她的监护责任撞在了一起。}她有没有权利未经你同意翻你手机？你说"没有"——她说"你未成年，我有监护义务"。\texbf{这两个答案都没有百分之百的错——它们在不同的"对"之间僵持。}

\subsection*{"这是我的群，我做主"}

班上有人建了一个微信群，专门用来截屏另一个同学的朋友圈，在里面嘲笑她穿的衣服、用的旧手机、发的"土味"自拍。群主被老师找去谈话。他双手一摊："这是我的群，我想拉谁拉谁，这是我的自由。"

他的自由——准确地停在什么位置？他确实有建群拉人的自由。但他使用了这项自由来把一个具体的、会哭的人当成了靶子。\textbf{自由给了你行使工具的权利——但没有给你用工具来撕碎别人的豁免权。}

\xuehai{密尔的"伤害原则"——自由的唯一合法边界}{
  19世纪的英国哲学家约翰·斯图尔特·密尔在《论自由》中提出了一个简洁到可以被刻在立法大楼门口的原则：\textbf{对任何一个文明社会成员，违背他意志地行使权力的唯一合法目的，是阻止他伤害他人。}他自己的好处——不管是身体上的还是道德上的——都不构成足够的理由。

  密尔的这句话翻译到你们班群的场景就是：你建群是你的自由。你拉谁是你的自由。你在群里说什么也是你的自由。但当群的唯一存在意义是\textbf{让一个特定的、有名有姓的人持续地被攻击时}——它不是在"行使自由"，它在用自由这件外套包裹一个伤害具体个体的武器。

  当然密尔的难题也在这里：\textbf{谁来定义"伤害"？}你朋友说"截我的图让我感到了伤害"——群主说"他自己心理太脆弱关我什么事"。"伤害"是一把必须由被伤害者来握住的尺子——密尔给了尺子，没给你握住任何一个具体人的手。
}

\subsection*{安全——自由的代价}

飞机安检让你把所有金属物品和液体丢进篮子里、笔记本电脑从包里拿出来、外套脱掉。你觉得很烦。为什么安检有权力碰你的私人物品？因为飞在万米高空的客舱是一个非常特殊的密闭空间——一个错误就可以让飞机上所有人付出生命的代价。你选择要坐飞机——等同于接受为了所有人的安全，被检查几秒钟行李。\textbf{自由和安全之间有一根谈判线}——永远在来回调整。

\subsection*{想一想}

\begin{enumerate}
  \item 你妈翻你手机——如果你是你妈，你会怎么做？有没有一种方式可以同时保护孩子的安全和尊重孩子的边界——还是必须二选一？
  \item 你的朋友每天截图你的朋友圈发给别人——这没犯法。但他的自由和你不愿被一直窥探的自由——哪一头在你心里更重？
  \item 密尔说"阻止他伤害别人"是唯一的干涉理由。如果一个人每天晚上在自家院子里对着墙反复喊同一句粗话——墙外没人听到——邻居也没有抱怨。你作为一个"觉得这种行为不好"的旁观者有权利去干涉他吗？密尔会怎么说？
\end{enumerate}

\newpage
\section{\textcolor{orange!80!black}{第六课\quad 法律——道德和法律的边界在哪？}}

\subsection*{邻居半夜在打鼓}

假设你家隔壁搬来的人喜欢凌晨两点练架子鼓。你严重睡眠不足、期末考试的复习计划被碾碎。你去找他说了两三次没效果。这时候你求助于——\textbf{法律}。有一条城市噪音管制条例明确写着：居住区夜间超过多少分贝、何时不得进行能穿透墙壁的高噪音作业。你报警——警察敲开他的门。

\subsection*{一部法律怎么从纸面掉到地上——跟着一条真实条款走一遍}

2021年6月，新修订的《\textbf{未成年人保护法}》开始施行。这次修订里加了一条在很多家长群和班级群里引起话题的条款：\textbf{学校、幼儿园周边不得设置烟、酒、彩票销售网点。}

这不是"写上去就有人自动遵守"。一个城市里开了十几年的学校门口小卖部——老板的柜台玻璃下面压着两三种不同价位的烟——每天早上有成年人靠在玻璃柜上买烟。新法生效之后，\textbf{市场监管局的执法人员按规定可以上门检查}。第一次检查——口头警告、下达整改通知书、限期腾空烟品货架。如果在后续的不定期回访中那条玻璃柜下仍然放着烟——店里可能被没收违规售卖的烟品并被处以几百至上万元不等的行政罚款。

法律从一个法条的字节变成了一纸处罚通知书再变成了小卖部玻璃柜台下那个\u200c\textbf{空了的格子}。这就是法律"走路"的全部动力。

\subsection*{撒谎、虚假和法庭}

你跟同学说"我没看那集动画"——其实你看了。撒谎了。道德上——你不诚实。但法律不会进你家来逮捕你，因为你并没有对着法官或公权力的文件撒谎。可是如果你在合同上签字谎报货款、或者在法庭上做了伪证——那一句假话可以让你坐牢。\textbf{道德的管辖领域远比法律大；但法律用于一旦引爆会产生严重社会伤害的行为——而且要动用公权力执行。}

\xuehai{法治和"用法来治"——一字之差的根本区别}{
  英国宪法学者戴西在1885年总结了"法治"的三个核心要素：第一，\textbf{没有人可以不经法律审判被惩罚}——惩罚你之前必须走过一遍公开的法律程序；第二，\textbf{所有人——包括政府官员——都服从同一部法律}（不算"总理的车可以闯红灯"）；第三，\textbf{宪法权利不是一纸宣言，而是法院可以在实际案件中强制执行的具体规则}。

  把这三条反过来就变成了"用法来治"：统治者用法律只约束别人不约束自己、法律惩罚没有公开程序、法院不能独立审判政府官员。同样的文字工具——"法律"——握在不同的人手里可以是盾也可以是矛。区别不在条文在页面上——在于\textbf{制定法律的人自己是否也被同一把尺子量}。
}

\subsection*{想一想}

\begin{enumerate}
  \item 学校门口的小卖部把烟收进了后台抽屉不摆出来——但老熟人问他还是偷偷从抽屉里拿一条——这算不算"遵守了法律"？法律的文字和法律的实质是同一个东西吗？
  \item 靠右行走——这是法律吗？遵守靠右行走到底是为了不被撞，还是因为这是人人默认的规矩？和你"不偷别人的笔"有什么本质区别？
  \item 戴西说"法治"意味着制定法律的人自己也被同一把尺子量。你们班的班规——班主任自己遵守吗？如果你发现一条规则只适用于学生但从不适用于发布这条规则的人——它更像"法治"还是更像"用法来治"？
\end{enumerate}

\newpage
\section{\textcolor{orange!80!black}{第七课\quad 我们和我们——个体、家庭、社区与国家}}

\subsection*{一个人可以独自做什么？}

你试着列一个清单——哪些东西你\textbf{完全不依赖任何别人}就能做成。从穿衣（衣服不是你织的）、吃的早饭（水电供应不是你自己拉的管道）到你正在看的这本书（有人排版有人写）。人类在几万年前走出了狩猎—采集小组之后，每跨一步需要的都是\textbf{超出单个身形的协作}。

\subsection*{一块荒地怎么变成了小广场}

你所在的城市有个老旧小区。两栋多层楼之间夹着一块空置了很多年的地。生锈的破沙发、碎玻璃渣、野猫窝。夏天有蚊虫、冬天灰扑扑。物业打扫过一次就不管了。

有一天，三楼的张阿姨整理了一颗花种子种在靠墙的一个旧轮胎里埋了土。六楼的李爷爷把他家的水管拿下来接了段皮管往轮胎里浇水。对面楼的妞妞放学看见有个轮胎开了一朵蓝紫色喇叭花，她回家跟自己爸爸说"我也想去种一颗"。几个月里陆陆续续有七八户人把破沙发清走了、铺了几十块捡来的废红砖。没有人发通知——大家你搬一块砖我拔一把野草——到第二年春天荒地的地面已经成了一块可以摆上塑料凳聊天的小坝子。

这块地\textbf{不是政府修的}。没有上头的什么"专项资金"，也不需要申请批文。它就靠着邻居之间自发协作——从一个所有人都绕道走的垃圾死角，变成了一块夏天有人在砖地上坐着扇扇子的地方。\textbf{有些事政府做不了——因为它看不见每一个小区中间的缝隙。但是看得见缝隙的人可以做。}

\subsection*{为什么一定要有国家？}

反过来——建一条连接两个城市之间的高速铁路。你自己挖土没有挖掘车，没有人把沿线一千户人的房子产权理清，没有统一的控制中心调度几百列车的速度——\textbf{这件事个人绝对做不了，市场做也会半路中断}。国家在这时候进来调动跨行政区、跨行业的生产链条。

\textbf{个体能做小事，社区能做中事，国家能做个人和社区都做不了的大事——这不是谁比谁更厉害，是工具和体量不同。}

\xuehai{邓巴数——为什么超过150人就需要规则}{
  英国人类学家罗宾·邓巴在研究灵长类动物大脑皮层尺寸和群体规模之间的关系时，发现了一个被反复验证的数字：\textbf{人类能够维持稳定社会关系的人数上限大约是150}。超过这个数——你没办法靠"我认识这个人、我信得过他"来协调行为——你需要规则、层级和制度。

  一个班级50个人——班主任能叫出所有名字。一个学校2000人——校长不可能认识每一个学生，需要年级组长、班主任、班委会逐层传递信息。一个城市两千万人——没有红绿灯没有警察局——靠邻里互相认识来维持秩序已经完全失效。邓巴数给"为什么小社区可以靠自觉、而大城市必须靠制度"提供了一道认知科学的地基。
}

\subsection*{被限制的权力}

国家自己可不可以想干嘛干嘛？\textbf{每一个现代的宪法都在回答这同一件事：不可以。}司法要独立于政治老板的喜好——否则法官看见某个人的背景就不敢下判决。立法的人不应同时执法——否则可以给你立个法律然后自己来抓你。\textbf{所有组织共同存在于国家的核心前提就是——国家本身的权力要被锁在一个笼子里。}

\subsection*{想一想}

\begin{enumerate}
  \item 小区那几位大爷阿姨把荒地变成了广场——你住的地方有类似的"没人管的角落"吗？你想怎么改？
  \item 什么事是你自己一个人绝对做不了的？（找出至少三件，看看它们分别需要哪一级力量来解决）
\end{enumerate}

\vspace{0.5cm}
\begin{center}
\textcolor{blue!70!black}{\large —— 从操场上的规则到广场石板上刻的法条，从一块荒地到一部宪法——道德和法治的拼图，你已经一块块拼完了。——}
\end{center}

\end{document}
